\chapter{Conclusion}
\label{ch:conclusion}

This thesis asked whether monocular depth estimation is good enough to replace a dedicated range
sensor for robot localization (Section~\ref{sec:research-question}). The Cyclops rover showed that
it is. With the depth calibration and range scan described in Chapter~\ref{ch:mde}, a single camera
gave Monte Carlo localization enough information to track the rover's pose in a known map and to
drive the rover to a goal, reaching the goal in $10$ of $10$ trials. This answer holds under the conditions listed in
Section~\ref{sec:conclusion-conditions}, which follow from how the depth is calibrated and from
the short range of the scan.

\section{Summary of Findings}
\label{sec:conclusion-findings}

The three parts of the research question were answered as follows.

\textbf{Metric depth without an additional sensor.} Depth Anything V2 predicts relative disparity,
which is correct only up to an unknown scale and shift that change from frame to frame
(Section~\ref{sec:alignment}). Because the camera is rigidly mounted above a flat floor, the depth
of the floor directly in front of the rover is known and fixed. Fitting the scale and shift to
these ground points on every frame, by ridge regression, recovered metric depth with no sensor
other than the camera (Section~\ref{sec:metric-scale}). The disparities of the ground points lay
close to a single line (Figure~\ref{fig:disparity-fit}), which confirmed that the affine model of
the disparity holds on the rover.

\textbf{A range scan for a standard measurement model.} The metric depth map was back-projected
into a 3D point cloud, stripped of the floor, and reduced to a range scan of $16$ rays
(Section~\ref{sec:mde-to-range}). Because this scan has the same form as the scan of a planar
LiDAR, MCL used it with the beam model of Thrun et al.~\cite{thrun2005probabilistic} without any
change to the model itself. Only the parameters of the model and the geometry of the ray casting
were adapted to the camera (Section~\ref{sec:measurement-model}).

% \textbf{Localization and navigation.} With the MDE range scan as its measurement, MCL estimated the
% rover's pose to within \red{[X]}~cm in position and \red{[Y]}$^\circ$ in heading of the true pose
% recorded by the overhead camera (Section~\ref{sec:results-localization}), and the rover reached its
% goal in \red{[N]} of $10$ trials (Section~\ref{sec:results-navigation}).
\textbf{Localization and navigation.} With the MDE range scan as its measurement, MCL localized the
rover well enough for it to reach its goal in $10$ of $10$ trials
(Section~\ref{sec:results-navigation}).

\section{Conditions and Limitations}
\label{sec:conclusion-conditions}

The answer above holds under the conditions that the method relies on.
\begin{itemize}
  \item \textbf{Lighting.} MDE predicts depth from the appearance of the image, so the rover must
  operate in a well-lit area (Chapter~\ref{ch:mde}).

  \item \textbf{A flat floor.} Metric scale is recovered from the floor, so the rover must drive on
  a flat ground plane, and anything that rises above that plane is treated as an obstacle
  (Section~\ref{sec:metric-scale}).

  \item \textbf{A clear strip in front of the rover.} The ground points used for calibration must
  see bare floor, so the rover cannot drive close enough to an obstacle for the obstacle to cover
  them, and a frame with no usable ground points returns no depth. This sets a minimum distance
  between the rover and any obstacle. In small scenes, it left few places where the rover could
  be, so the tests were run in a more spacious scene (Section~\ref{sec:results-limitations}).

  \item \textbf{Short range.} Ranges were used only up to $0.40$~m, so the rover saw only the
  obstacles close in front of it.

  \item \textbf{A known map and a small test area.} The rover localized against a map made in
  advance, and all tests were run inside a $1.105$~m $\times$ $1.82$~m enclosure with one
  arrangement of obstacles (Section~\ref{sec:test-env}).

  \item \textbf{Off-board computation.} Depth Anything V2 ran on a laptop rather than on the rover,
  because it was too slow on the Raspberry Pi (Section~\ref{sec:DAV2}). At about three range scans
  per second, the rover had to move in short steps and stop to wait for a new scan after each one
  (Section~\ref{sec:navigation}).
\end{itemize}

Chapter~\ref{ch:future-work} describes ways to relax several of these conditions, such as
computing the ground points from the camera's mounting geometry instead of a checkerboard, and
learning the parameters of the measurement model from data.

Within these conditions, the Cyclops rover showed that a single camera can do the work of a
dedicated range sensor for localization. A small robot that already carries a camera can
therefore localize itself in a known map without the added cost, weight, and power draw of a
LiDAR.

%%% Local Variables: ***
%%% mode: latex ***
%%% TeX-master: "thesis.tex" ***
%%% End: ***
